% !TeX root = ../../../manual.tex
\chapter{Floats}
\label{ch:floats}

\section{Float Placement}

KOMA-Script extends the standard \LaTeX{} float placement with additional
specifiers. OmniLaTeX loads the \texttt{floatrow} package to provide
automatic caption positioning and improved float layouts.

\begin{longtblr}[
    caption = {Float placement specifiers},
    label = {tab:float-placement},
]{colspec = {c l}, rowhead = 1}
    \toprule
    Specifier & Meaning \\
    \midrule
    h & Place here if possible \\
    t & Place at the top of a page \\
    b & Place at the bottom of a page \\
    p & Place on a dedicated float page \\
    H & Place exactly here (requires \texttt{float} package) \\
    ! & Override internal \LaTeX{} placement constraints \\
    + & Allow floats to appear in the current region (KOMA extension) \\
    \bottomrule
\end{longtblr}

The default placement in OmniLaTeX is \texttt{htbp}. Override it globally or
per-float:

\begin{verbatim}
\begin{figure}[!htb]
    ...
\end{figure}
\end{verbatim}

\section{Continued Floats}

For figures that span multiple pages, use \verb|\ContinuedFloat| from the
\texttt{caption} package. The caption counter is not incremented and a
\emph{(continued)} marker is appended:

\begin{verbatim}
\begin{figure}[htbp]
    \caption{First part of a multi-page figure}
    \label{fig:long-first}
    \includegraphics[width=\textwidth]{part1}
\end{figure}

\begin{figure}[htbp]
    \ContinuedFloat
    \caption{Second part of a multi-page figure}
    \label{fig:long-second}
    \includegraphics[width=\textwidth]{part2}
\end{figure}
\end{verbatim}

\section{Float Footnotes}

Standard \LaTeX{} footnotes inside floats are fragile and may not survive
page breaks. OmniLaTeX provides a robust mechanism:

\begin{verbatim}
\omnlFloatFootmark{fn1}
...
\omnlFloatFoottext{fn1}{This footnote survives page breaks.}
\end{verbatim}

The footnote mark is rendered inline at the point of invocation, while the
footnote text is placed below the caption. Multiple keys can be used within
a single float.

\section{Side Captions}

Place captions beside figures using a \texttt{minipage} layout:

\begin{verbatim}
\begin{figure}[htbp]
    \begin{minipage}[c]{0.65\textwidth}
        \includegraphics[width=\textwidth]{photo}
    \end{minipage}%
    \hfill
    \begin{minipage}[c]{0.30\textwidth}
        \captionof{figure}{Description alongside the image.}
        \label{fig:side-caption}
    \end{minipage}
\end{figure}
\end{verbatim}

\section{Key-Value Float Interface}

OmniLaTeX provides high-level commands for figures and tables using key-value
syntax:

\begin{verbatim}
\omnlFigure[
    caption        = {An example figure},
    label          = {fig:kv-example},
    placement      = htbp,
    caption-width  = 0.8\textwidth,
    caption-position = bottom,
]{%
    \includegraphics[width=\textwidth]{example-image}
}
\end{verbatim}

\begin{longtblr}[
    caption = {\texttt{\textbackslash omnlFigure} options},
    label = {tab:omnlfigure-options},
]{colspec = {l l}, rowhead = 1}
    \toprule
    Key              & Description \\
    \midrule
    caption          & Caption text \\
    label            & Label for cross-references \\
    placement        & Float placement specifiers (default: \texttt{htbp}) \\
    caption-width    & Width of the caption box \\
    caption-position & \texttt{top} or \texttt{bottom} (default: \texttt{bottom}) \\
    source           & Source attribution text \\
    \bottomrule
\end{longtblr}

The \verb|\omnlTable| command accepts the same keys and wraps content in a
\texttt{table} environment.

\section{Float Grid}

For multi-column float layouts, use \verb|omnlFloatRow|:

\begin{verbatim}
\omnlFloatRow[n-cols=3]{%
    \includegraphics[width=\linewidth]{img1}%
    \includegraphics[width=\linewidth]{img2}%
    \includegraphics[width=\linewidth]{img3}%
}
\end{verbatim}

This creates an evenly spaced row of images. The optional \texttt{n-cols}
parameter defaults to 2.

\section{Figure Environment}

The standard \texttt{figure} environment is fully supported with OmniLaTeX
enhancements:

\begin{verbatim}
\begin{figure}[htbp]
    \centering
    \includegraphics[width=0.9\textwidth]{diagram}
    \caption{A standard figure with OmniLaTeX styling.}
    \label{fig:standard}
\end{figure}
\end{verbatim}

OmniLaTeX automatically applies the configured caption font and separator
style to all float captions.

\section{Table Environment}

Similarly, the \texttt{table} environment benefits from OmniLaTeX defaults:

\begin{verbatim}
\begin{table}[htbp]
    \centering
    \begin{tabular}{lcc}
        \toprule
        Parameter & Value & Unit \\
        \midrule
        Mass      & 1.23  & kg  \\
        Length    & 4.56  & m   \\
        \bottomrule
    \end{tabular}
    \caption{A standard table with booktabs styling.}
    \label{tab:standard}
\end{table}
\end{verbatim}
